\section{Spatial initial metrics and the initial layer}\label{sec:entry}

The initial-error hypothesis of Theorem~\ref{thm:tracking}
can be obtained from spatially varying metrics on the
normalization components. We first select the constant
metric determined by the actual flow on each fixed
component. We then construct an initial layer on the
smooth fibre whose duration is negligible compared
with $L_*$.

\subsection{Constant limits on a fixed component}
Let $Y$ be a fixed connected compact Riemann surface
of area $A$, with its trivial holomorphic rank-$r$
bundle and fixed flat frame. Write
$\langle f\rangle=A^{-1}\int_Y f\,dA$, entrywise
for a matrix. Let $\lambda_1>0$ be the first
positive scalar eigenvalue.
When a component flow converges to a constant
metric $P$, we call $P$ the asymptotic phase
of its initial metric.

\begin{theorem}\label{thm:component-phase}
For a smooth determinant-one initial metric with
$cI\le H^{\mathrm{in}}\le CI$, $0<c\le1\le C$,
the actual HYM flow converges to a unique constant
matrix $P=P(H^{\mathrm{in}})\in\cP_r$.
Set
\[
 M(s)=\langle H(s)\rangle,\quad J(s)=\log\det M(s),
 \quad V(s)=\langle\|H(s)-M(s)\|_{\HS}^2\rangle,
 \quad \gamma=2c^2C^{-2}\lambda_1 .
\]
Then $M$ decreases in Loewner order to $P$, and
\begin{equation}\label{eq:phase-entropy}
 \frac{V}{2C^2}\le J\le\frac{V}{2c^2},\quad
 J(s)\le J(0)e^{-\gamma s},\quad
 0\preceq M(s)-P,\quad \tr(M(s)-P)\le CJ(s).
\end{equation}
In particular
\begin{equation}\label{eq:phase-l2}
 \langle\|H(s)-P\|_{\HS}^2\rangle
 \le C^2(2J(0)+J(0)^2)e^{-\gamma s}.
\end{equation}
For $s\ge1$ the same exponential order bounds
$\sup_Y d(H(s),P)^2$, with a constant depending
only on $Y,c,C,r$. For two such initial metrics,
\begin{equation}\label{eq:phase-contraction}
 A\,\sigma(P(H^{\mathrm{in}}),P(K^{\mathrm{in}}))
 \le\int_Y\sigma(H^{\mathrm{in}},K^{\mathrm{in}})\,dA .
\end{equation}
\end{theorem}
\begin{proof}
Equation~\eqref{eq:local-hym} is
$H_s=-\Delta H-\mathcal A(H)$, with
$\mathcal A(H)=4a^{-1}H_z^*H^{-1}H_z\succeq0$.
Its spatial mean therefore satisfies
$M'=-\langle\mathcal A(H)\rangle\preceq0$.
Since $cI\le M\le CI$, it has a positive matrix
limit $P$. At this point its determinant has
not been prescribed.

For a Hermitian direction $Z$ and a matrix
$cI\le B\le CI$,
\[
 -D^2(\log\det)_B[Z,Z]
 =\tr(B^{-1}ZB^{-1}Z),\qquad
 C^{-2}\|Z\|^2\le\tr(B^{-1}ZB^{-1}Z)
                  \le c^{-2}\|Z\|^2.
\]
Diagonalizing $B$ proves the inequalities
entrywise. Apply Taylor's integral remainder
along the segment from $M$ to $H(x)$, then
average. The linear term vanishes and
$\log\det H(x)=0$, giving the two bounds
for $J$ in \eqref{eq:phase-entropy}.

Using positivity twice gives
\[
 -J'=\tr(M^{-1}\langle\mathcal A(H)\rangle)
 \ge C^{-1}\langle\tr\mathcal A(H)\rangle
 \ge C^{-2}\langle|dH|^2\rangle.
\]
In the last step one first bounds $H^{-1}$
below and then uses the unweighted identity
$4a^{-1}\tr(H_z^*H_z)=|dH|^2$. The scalar
Poincar\'e inequality in each matrix coordinate
gives $\langle|dH|^2\rangle\ge\lambda_1V$.
Thus $J'\le-\gamma J$. In particular
$J\to0$, so continuity of determinant gives
$\det P=1$.

The segment from $P$ to $M$ is bounded above
by $CI$. Since $M-P\succeq0$,
\[
 J=\int_0^1\tr\!\left[
 (P+t(M-P))^{-1}(M-P)\right]dt
 \ge C^{-1}\tr(M-P).
\]
The mean-zero cross term in the squared norm
vanishes, and therefore
\[
 \langle\|H-P\|^2\rangle
 =V+\|M-P\|^2\le2C^2J+C^2J^2.
\]
This proves \eqref{eq:phase-l2}.

The constant $P$ is itself an exact solution.
The scalar trace comparison shows that
$\sigma(H(s),P)$ is a nonnegative heat
subsolution. On the fixed component the
time-one heat kernel is bounded, whence
\[
 \sup_Y\sigma(H(s),P)
 \le C_Y\int_Y\sigma(H(s-1),P)\,dA
 \quad(s\ge1).
\]
On the fixed matrix range $\sigma$ is
comparable to squared matrix difference,
so \eqref{eq:phase-l2} proves uniform
convergence with the asserted rate.
Finally apply trace contraction to two
exact flows and pass to their uniform
constant limits. This gives
\eqref{eq:phase-contraction}.
\end{proof}

The phase depends on the entire spatial initial
metric. For diagonal initial data
$\operatorname{diag}(e^{f_1},\ldots,e^{f_r})$,
$\sum f_j=0$, it is
$\operatorname{diag}(e^{\langle f_1\rangle},
\ldots,e^{\langle f_r\rangle})$, by scalar
heat flow. No formula
$P=\exp\langle\log H^{\mathrm{in}}\rangle$
is asserted for general noncommuting data.
The contraction estimate makes the phase map
Lipschitz in $L^2$ on every fixed bounded
metric range.

\subsection{An open class on the smooth fibres}
Put $\rho=L_*^{-1}$ and let
$\Omega_i(2\rho)\subset C_i$ be the complement
of the incidence discs of radius $2\rho$.
Use the fixed exterior transfers and flat frames
of Section~\ref{sec:geometry}.
For a smooth trace-free $h_0$-Hermitian
endomorphism $S$ on $C_\tau$, fix $0<\beta<1$
and define
\begin{equation}\label{eq:extension-norm}
 \mathcal N_\tau(S)=\|S\|_{C^0(C_\tau)}
 +\inf_{\widetilde S_i|_{\Omega_i(2\rho)}=S_i}
             \sum_i\|\widetilde S_i\|_{C^{6,\beta}(C_i)} .
\end{equation}
The extensions are trace-free Hermitian matrices
in the component frames. Fix $R_0>0$ and take
$\mathcal N_\tau(S)<R_0$, with initial metric
$H^{\mathrm{in}}=h_0e^S$.
For each fixed smooth fibre this class is open
in the $C^{6,\beta}$ topology on smooth metrics:
fixed-domain extension operators bound
$\mathcal N_\tau(S_1-S_2)$ by a finite multiple
of the global $C^{6,\beta}$ norm.
Their operator norms need not be uniform as
$\rho\to0$. The radius in
\eqref{eq:extension-norm} is fixed independently
of the degeneration. Small fixed component
profiles joined smoothly on the nodes give
nonempty examples. Bounded smooth changes
inside the middle necks may have arbitrary
derivatives.

Choose extensions with total norm less than
$2R_0$. Apply Theorem~\ref{thm:component-phase}
to $e^{\widetilde S_i}$ and denote its flow by
$G_i(s,z)$ and selected phase by $P_i$.
All component flows have uniform positive
bounds and uniform exponential phase convergence.
Different admissible extensions agree outside
discs of area $O(\rho^2)$, so
\eqref{eq:phase-contraction} gives
$d(P_i,\widetilde P_i)=O(\rho)$.
Finite-flow nonexpansion preserves this
ambiguity at all later times.

\begin{theorem}\label{thm:entry}
For the class \eqref{eq:extension-norm} there
is a preparation time $s_*(\tau)\to\infty$,
uniform in all length ratios and unitary
parameters, such that $s_*/L_*\to0$.
Let $\widehat K_P(v)$ be the reconstruction
of Theorem~\ref{thm:tracking} for the finite
flow starting at the selected tuple $P$.
Then
\[
 \eta_*=\int\sigma(H(s_*),\widehat K_P(0))\to0
\]
uniformly over the stated initial class.
Consequently \eqref{eq:tracking} holds with
$H(s_*+v)$, $\widehat K_P(v)$ and initial
error $\eta_*$, for $0\le v\le TL_*$.
Under \eqref{eq:rate-condition} the all-time
estimate of Theorem~\ref{thm:all-time} holds
after the same preparation.
\end{theorem}
\begin{proof}
Fixed-base parabolic estimates for the component
flows with bounded $C^{6,\beta}$ initial profiles
give nondecreasing finite constants $C_n\ge1$
controlling the required space and time derivatives
on $0\le s\le n$. They depend on $R_0$ and the
fixed components, not on $\tau,U$.
Transfer $G_i(s,z)$ to the component exteriors.
On each node join its two point values
$G_i(s,p_{ie})$ and $U_{ij}G_j(s,p_{je})U_{ij}^*$
by the geodesic with parameter
$-\varepsilon_e\log|u_e|$. The node curvature
is zero, as in \eqref{eq:node-profile}.
On the $i$ collar use the spatially constant
reference $G_i(s,p_{ie})$ for both relative
logarithms. On $\rho\le|u_e|\le2\rho$, the difference
between the component and node logarithms, with
its first two derivatives, is at most
\[
 C_n(\rho+\varepsilon_e|\log\rho|)\rho^{-k},
 \qquad k=0,1,2 .
\]
The $O(\rho)$ term is the spatial variation
of $G_i$ from its nodal value, and the
other term is the small node parameter.
Use the same logarithmic cutoffs and
curvature-difference estimate as in
\eqref{eq:mixed-collar}. The integrated
collar residual is at most
$C_n(\rho+\varepsilon_*|\log\rho|)$.
The middle-node time residual is
$O(C_n\rho^2)$ and the geometric transfer
error is $O(C_nb\rho^{-N})$.
Thus the resulting global preparation
metric $K_{\mathrm{pre}}$ obeys
\[
 \sup_{0\le s\le n}\int|\mathcal R(K_{\mathrm{pre}})|
 \le C_n\delta_*,
 \qquad
 \delta_*=\frac{1+\log L_*}{L_*}+e^{-L_*}L_*^N\to0 .
\]
The initial metric and $K_{\mathrm{pre}}(0)$
agree outside regions of area $O(\rho^2)$
and have common positive bounds. By
\eqref{eq:trace-comparison} their integrated squared
distance, and their integrated trace
discrepancy, are bounded by
$C\rho^2+nC_n\delta_*$ up to time $n$.

For all sufficiently small parameters choose
\[
 s_*=\max\{n\in\mathbb N:nC_n\le\delta_*^{-1/2}\}.
\]
It tends to infinity because each $C_n$ is
finite, and it is finite because $C_n\ge1$.
Also $s_*\le\delta_*^{-1/2}$, which implies
$s_*/L_*\to0$. At that time the component
flows are exponentially close to $P_i$.
The reconstruction with initial tuple $P$
differs from the constant component metrics
by the Green corrections; their squared
integrals are $O(L_*^{-2})$. The remaining
node regions have area $O(\rho^2)$.
On the common matrix range, $\sigma$ is
comparable to squared distance and obeys
a triangle bound up to a fixed constant.
We obtain
\[
 \eta_*\le C_{R_0}\left(
 \sqrt{\delta_*}+e^{-\gamma_{R_0}(s_*-1)}
 +L_*^{-2}\right)\to0 .
\]
Restart Theorem~\ref{thm:tracking} at
$s_*$, and Theorem~\ref{thm:all-time} when
its additional rate hypothesis holds.
\end{proof}

For physical times after preparation one may align
the finite trajectory with time zero. Its speed
is $O(L_*^{-1})$, and the uniform $L^2$ Lipschitz
reconstruction gives, for $v\ge0$,
\[
 \int d(H(s_*+v),\widehat K_P(s_*+v))^2
 \le 2\int d(H(s_*+v),\widehat K_P(v))^2
       +C(s_*/L_*)^2 .
\]
Thus the time shift costs a fixed factor in the
preceding squared-error estimate and a vanishing
additional term. Uniform tracking does not
by itself assert convergence of a family of
selected phases. To pass to a fixed positive
rescaled-time limit, one must also assume
convergence of those phases and of the
normalized finite parameters.
